% ARQUITECTURA_AUTORAL_PREEXISTENTE: bidirectional half-square of N42.
% FORMALIZACION_NUEVA: pair groupoid and stabilizer weighting.
\subsubsection{Marked incidence and bidirectional normalization}
\label{sec:ii09-normalizacion-bidireccional}

The reader invariants preserve the nature of the object from which
they arise. The marked intersection $B_K=H_e\cap P_\pi$ has cardinality
$b=4$; the integer $q=7$ is the coordinate of the pivot $p_\varphi$ in
the oriented numbering. The determination order of the Steiner
design is $p=5$. The hexadic and octadic supports have cardinalities $h_6=6$
and $o_8=8$, and the tritic orientation alphabet has cardinality $t=3$.
A support cardinality and a pivot coordinate are used through
their respective maps, even when both are expressed as integers.

Normalization of the pentagonal self-crossing admits the following
explicit realization. Let $P$ be the five-position support of the closure,
and let $X=P\times P$ be its space of ordered pairs. Bidirectional
inversion on this space is
\[
 \jmath:X\longrightarrow X,\qquad \jmath(x,y)=(y,x),
 \qquad\jmath^2=I.
\]
The stabilizer information of each orbit is preserved. For an
orbit $[z]$, define its weight as
$1/|\operatorname{Stab}_{C_2}(z)|$. The sum of these weights is the
cardinality of the action groupoid $X/\!/C_2$.

\begin{proposition}[Orbital measure of bidirectional self-crossing]
\label{prop:ii09-semicuadrado-orbital}
For a finite support of cardinality $p$, inversion of ordered pairs
determines
\[
 |(P\times P)/\!/C_2|
 =\sum_{[z]\in(P\times P)/C_2}
      \frac1{|\operatorname{Stab}_{C_2}(z)|}
 =\frac{p^2}{2}.
\]
For $p=5$, the self-crossing has ten free orbits and five fixed
orbits, and its orbital measure is $10+5/2=25/2$.
\end{proposition}
\begin{proof}
Each diagonal pair $(x,x)$ forms a fixed orbit, whose stabilizer
has order two; there are $p$ such orbits. The remaining $p(p-1)$ pairs
are grouped in twos, with trivial stabilizer. Their contribution is
therefore $p(p-1)/2$, and that of the diagonal is $p/2$. The sum is
$p^2/2$. Equivalently, the orbit--stabilizer identity gives
$1/|\operatorname{Stab}(z)|=|[z]|/2$; summing over all orbits
produces $|P\times P|/2$.
\end{proof}

The measure explicitly preserves the contribution of fixed points.
The ordinary orbit set has cardinality $p(p+1)/2$, equal to
fifteen in the pentagonal closure. The orbital measure and that cardinality
are two determinate readings of the same inversion. The reader's
bidirectional normalization uses the former: the ten free orbits
have unit weight, and the five diagonals have half-unit weight.

This realization specifies the factor $p^2/2$ used in the following
sectorial system. The assignment of that contribution to sector $23$,
its negative orientation, and the other two incidence rules remain
specified by the complete sectorial system. The proof establishes the
normalization of self-crossing and preserves the selection of the
angular functionals as a distinct matter.
